\section{Implementació de DDPG i TD3}\label{sec:impl}

DDPG \parencite{lillicrap2016} aprèn un actor determinista $\mu(s)$ i un crític $Q(s,a)$ a partir d'una memòria de repetició. El crític minimitza l'error quadràtic respecte de l'objectiu $y = r + \gamma(1-d)\,Q'(s',\mu'(s'))$, calculat amb les xarxes objectiu $Q'$ i $\mu'$, i l'actor minimitza $-Q(s,\mu(s))$. Després de cada iteració, les xarxes objectiu s'actualitzen suaument amb el coeficient $\tau = 0{,}001$. A $d$ només hi compta la terminació real (estavellar-se, sortir de la pantalla o quedar en repòs): quan l'episodi es talla pel límit de 1.000 passos, l'objectiu continua fent servir el valor de $s'$, perquè aquest estat no és terminal. Es fa una iteració per pas d'entorn.

Els hiperparàmetres són els de la secció 7 de \textcite{lillicrap2016}, com demana l'enunciat: Adam amb taxa d'aprenentatge $10^{-4}$ a l'actor i $10^{-3}$ al crític, \emph{weight decay} de $10^{-2}$ al crític, lots de 64, memòria de $10^6$ transicions i factor de descompte $\gamma = 0{,}99$. Les xarxes són les de l'esquelet, sense modificar (400 i 300 unitats amb LayerNorm). El constructor del crític no crida la inicialització de l'article i la cridem explícitament. DDPG explora amb el soroll d'Ornstein-Uhlenbeck (OU) de \texttt{Noise.py} amb els paràmetres per defecte ($\theta = 0{,}2$, $\sigma = 0{,}15$, $dt = 0{,}01$; l'article usa $\theta = 0{,}15$ i $\sigma = 0{,}2$), reiniciat a cada episodi.

TD3 \parencite{fujimoto2018} és una còpia de DDPG amb tres canvis. (i) Suavitzat de la política objectiu: $\tilde a' = \mathrm{clip}(\mu'(s') + \mathrm{clip}(\epsilon,\;{-0{,}5},\;0{,}5),\;{-1},\;1)$ amb $\epsilon \sim \mathcal{N}(0,\;0{,}2^2)$. (ii) Dos crítics entrenats cap al mateix objectiu $y = r + \gamma(1-d)\min_{i=1,2} Q'_i(s',\tilde a')$; l'actor maximitza $Q_1$. (iii) L'actor i les tres xarxes objectiu s'actualitzen cada dues iteracions dels crítics. Es mantenen $\tau$, les xarxes i els optimitzadors, inclòs el \emph{weight decay} dels dos crítics. L'exploració és gaussiana, $\mathcal{N}(0,\;0{,}1^2)$, el valor de l'article de TD3, perquè l'enunciat no en fixa cap.

Tres interruptors (\texttt{use\_minimum}, \texttt{target\_noise} i \texttt{policy\_delay}) defineixen les variants de l'ablació sense canviar res més. \emph{Sense mínim} fa servir $Q'_1$ a l'objectiu però continua entrenant els dos crítics, \emph{sense suavitzat} posa $\sigma_\text{targ} = 0$ i \emph{sense retard} actualitza l'actor i les xarxes objectiu a cada iteració. El soroll de l'objectiu se sorteja sempre, també amb $\sigma_\text{targ} = 0$: així, amb la mateixa llavor, totes les variants parteixen dels mateixos pesos i consumeixen la mateixa seqüència aleatòria. Proves automàtiques comproven aquestes propietats.
